% SPDX-FileCopyrightText: 2026 David Shirvanyants
% SPDX-License-Identifier: CC0-1.0

\chapter{Interface Manual}
\label{interface-manual}

\section{Document Window}
\label{document-window}

Each document contains a model list, a toolbar, and an OpenGL viewing area.
The document window is the normal workspace for one or more loaded models.
Use the \hyperref[model-list]{model list} to select models and control their
visibility, the \hyperref[document-toolbar]{toolbar} or
\hyperref[document-menus]{menus} to start operations, and the viewport to
inspect or transform geometry. Detailed gesture descriptions are collected
under \hyperref[mouse-actions]{Mouse Actions}; keyboard access is summarized
under \hyperref[keyboard-shortcuts]{Keyboard Shortcuts}.
On Linux, each document tab has a close button beside its title. On Windows,
documents retain the native MDI window controls supplied by the operating
system.

Open and export file dialogs retain the directory of the last accepted file
for use by subsequent dialogs. The model viewport requires OpenGL 3.2 and GLSL
1.50. If renderer initialization fails, the application reports the detected
OpenGL environment and a sample of missing entry points to assist driver
diagnosis.

\subsection{Model List}
\label{model-list}

The model list controls selection and visibility. Its check boxes show or hide
model geometry without deleting it. Commands and transformations apply to all
selected entries, and a document can contain both source surface models and
generated slice models.

\section{Menu Commands}
\label{document-menus}

The \emph{File} menu opens a model in a new document or in the active document,
exports selected slices as CLI, exports selected surface models as STL, opens
Settings, and closes the document or application.
\hypertarget{open-model-dialog}{}The \emph{Open} command uses the Open Model
dialog and creates a new document.
\hypertarget{open-into-document-dialog}{}The \emph{Open into document} dialog
adds the chosen STL or CLI file to the active document.
\hypertarget{export-slices-dialog}{}The \emph{Export Slices} dialog writes the
selected sliced models to a CLI file.
\hypertarget{export-stl-dialog}{}The \emph{Export STL} dialog writes selected
surface models to a binary STL file.

The \emph{Models} menu contains commands that act on every selected model.
\emph{Delete selected} removes all selected models from the current document.
\emph{Reset transformations} restores each selected model's original position,
orientation, and scale. \emph{Transform} combines exact rotation, translation,
and uniform scaling in one dialog. Rotation uses a chosen global X, Y, or Z
axis; rotation and scaling act on the selected models as a group around their
common geometric center, followed by signed translation along the three global
axes. These operations update transformation matrices without changing source
model coordinates. \emph{Move to Origin} preserves the relative placement of
all selected models while translating their combined bounding box so that its
minimum X, Y, and Z coordinates are all zero. The selected group therefore
fits tightly in the positive octant and rests on the build plane.

\hypertarget{transform-dialog}{}The \emph{Transform} dialog provides numeric
rotation axis and angle, translation, and uniform scale controls.

The \emph{Slice} menu contains commands for regular slicing, interactive
cross-sections, unsupported-area detection, orientation optimization, and
support generation. \hypertarget{slice-dialog}{}The \emph{Slice selected}
dialog chooses whether generated slices remain in the same document or are
placed in a new document. The \emph{Help} menu's \emph{Topics} command opens
the top of this manual.

\section{Toolbar Buttons}
\label{document-toolbar}

The toolbar provides the most frequently used menu commands. \emph{Export},
\emph{Slice}, and \emph{Section} write or inspect layers. \emph{Hide},
\emph{Show}, and \emph{Delete} affect selected list entries. \emph{Detect},
\emph{Optimize}, and \emph{Generate} run the analysis and support workflows;
\emph{Stop} cancels a running background workflow. \emph{Reset},
\emph{Transform}, and \emph{Origin} change selected-model transformations.
Buttons which cannot operate on the current selection are disabled.

\section{Viewport Guides}
\label{viewport}

\subsection{Build Platform}

The rectangle below the models represents the printer build platform. Its
width and length follow the combined transformed bounds of all visible models.
Its upper surface lies on the global Z=0 plane.
The upper surface is dark gray. The underside remains visible when viewed from
below and is shown in silver gray, making the platform orientation apparent.
The first section of the application status bar reports the current build
volume as width, length, and height.

\subsection{Coordinate Axes}

The three thin lines crossing the viewport are the global X, Y, and Z axes.
They pass through the global origin and extend in both positive and negative
directions until clipped by the view. Cross-shaped tick marks remain visible
from different viewing directions. Tick spacing changes with zoom and always
uses a power of ten, such as 0.1, 1, 10, or 100 model units.

The axes belong to the global coordinate system. Camera movement changes their
screen projection, while model transformations move models relative to them.

\subsection{Orientation Vane}

The orientation vane is the group of three thick arrows in the lower-left
corner of the viewport. Its X arrow is blue, its Y arrow is red, and its Z
arrow is bright green. Black X, Y, and Z labels identify the corresponding
arrow ends. The vane rotates with the camera to show the current view
direction, but its screen position and size remain fixed so that it is always
legible.

\subsection{Cross-section}
\label{section-dialog}

The \emph{Section} toolbar button or \emph{Slice > Section} menu command opens
a dialog for choosing an X-, Y-, or Z-normal section plane and optionally
aligning the camera for the best view. Its Clipping choice can remain \emph{No},
or hide selected model geometry \emph{Above} or \emph{Below} the moving plane.
The translucent section plane spans the
shared bounds of the selected models and can be moved along its normal with
Alt+wheel. Plane positions follow the configured slicing sequence: the first
position for the Z axis is aligned to the global build-layer grid defined by
the first-layer offset and layer thickness. X- and Y-axis sections start at
the selected-bounds minimum plus the first-layer offset. Each subsequent
position advances by one layer thickness. While Section mode is
active, a horizontal scrollbar above the view and a spin box to its right show
and control the slice index. Moving the thumb right increases the index, and
the projected cross-section updates continuously while the thumb is dragged.
With the scrollbar focused, the Up and Down keys move by one slice
and Page Up and Page Down move by ten slices.

\hypertarget{section-controls}{}These two controls constitute the Section
position controls and share the exact slice-index limits used by the plane.

Moving the plane recalculates the intersection contours for selected models.
A filled, double-sided copy of those contours is drawn on a dedicated
semi-transparent plane outside the selected models, and the second status-bar
field reports the coordinate using the selected axis letter. Best-view
alignment faces the display plane squarely and scales it to fit the viewport.
Hiding models affects their rendered geometry only: the build platform,
section and projection planes, and projected contours remain visible, including
when all models are hidden.

\subsection{Unsupported Areas}

The \emph{Detect} toolbar button or \emph{Slice > Detect unsupported areas}
command analyzes the selected models in their transformed positions. The
adjacent \emph{Generate} button and \emph{Slice > Generate supports} command
are the entry point for creating support structures.
The analysis uses the configured slicing thickness, critical support angle,
and overhang coefficient. It runs in the background so that document viewing
remains responsive.

Support generation snapshots the selected transformed models, slices them,
and detects unsupported areas on a background coordinator thread. Its complete
processing sequence is described in Chapter~\ref{chap:generate}.

When analysis finishes, unsupported portions of layer surfaces are overlaid in
bright red-orange. For clearer visibility against the source mesh, each
unsupported area is drawn at the Z position of the preceding layer rather
than the slicing plane on which it was detected. The first layer is treated
as fully supported. The layer-to-layer support test and the treatment of
disconnected islands are described in Chapter~\ref{chap:detect}. Transforming
a model clears the overlay because the previous result no longer represents
its current position.

\subsection{Orientation Optimization}

The \emph{Optimize} toolbar button or \emph{Slice > Optimize orientation}
command searches for orientations which reduce the total unsupported area of
the selected models. Selected models are optimized independently and in list
order. Improved orientations are applied to the document as they are found;
the source vertex coordinates remain unchanged.

The search procedure and its unsupported-area score are described in
Chapter~\ref{chap:optimize}. Several attempts run concurrently. The third
status-bar field reports
completed attempts and the best unsupported-area score for the current model
in the form \texttt{Run 1 of 10 (S=12.34)}. The current orientation is analyzed
first, and its initial score is displayed before rotated search candidates are
evaluated. Model transformations or analysis-setting changes cancel stale
optimization work.
During support generation, the same status field follows an ordered step list:
\emph{Slicing}, \emph{Finding unsupported}, \emph{Contact points},
\emph{Volume segmentation}, and \emph{Generating supports}. Contact discovery
and volume segmentation may run concurrently, but the display advances only
after the last worker for the current step has finished. Contact stages show
calculated or processed contacts in the form \texttt{(10 of 1234)}. Before the
final contact count is known, the total is estimated from unsupported area and
the configured contact spacing.
After generation, contacts for which neither an external nor internal support
could be placed are summarized in the same field. The message includes the
number of failures and the Z coordinate of the lowest failed upper contact.
The \emph{Stop} toolbar button requests cooperative cancellation. Detection,
optimization, and support-generation workers stop between slicing, analysis,
or generation operations, and the controls return to their idle state after
all workers have exited. Interactive Section previews are calculated directly
for immediate feedback and do not enable the Stop button.

\subsection{STL Export}

The \emph{File > Export STL} command writes all selected surface models into
one binary STL file. Stored model transformations are applied to every vertex
and normal. Sliced models are omitted because their visualization mesh is not
treated as source STL geometry. Intersecting input meshes are written without
performing a geometric union.

\section{Settings}
\label{settings-dialog}

The \emph{File > Settings} command opens the application settings dialog.
Settings are grouped into tabs so that additional categories can be added as
the application grows.

\subsection{Slicing Tab}
\label{settings-slicing}

The \emph{Slicing} tab defines the default layer thickness and absolute
first-layer offset above the build platform. These values define the shared
layer grid used by model slicing, unsupported-area analysis,
orientation optimization, and interactive cross-sections. They default to 0.1 and
0.05 model units, respectively.

The tab also contains the contour healing threshold. After initial
segment connection, endpoints separated by no more than this distance can be
joined to close a contour. The default is 0.01 model units. The setting is
used by normal slicing, interactive cross-sections, and unsupported-area analysis.

The segmentation tolerance controls endpoint matching during initial slice
segment connection. Its default is 0.01 model units. It is shared by the
same three slicing workflows.

\subsection{Analysis Tab}
\label{settings-analysis}

The \emph{Analysis} tab controls unsupported-area detection. The critical
angle defaults to 30 degrees. The overhang coefficient defaults to 1 and
scales the direct layer-to-layer allowance. For a layer separation \(dZ\),
the support search radius is
\(dZ (c + 1 / \tan(\alpha))\), where \(c\) is the overhang coefficient and
\(\alpha\) is the critical angle. Increasing the coefficient permits more
horizontal displacement between successive layers before an area is marked
unsupported.
When unsupported-area detection takes longer than one second, the status bar
reports the completed slice and total slice count. The report is updated at
most once per second to avoid slowing the analysis with GUI traffic.

The optimization-attempt setting controls the number of local searches per
model and defaults to 10. The worker-thread setting controls both the maximum
number of optimization attempts evaluated concurrently and the support
generation worker pool. It defaults to 4.
Convergence tolerance is the minimum unsupported-area reduction accepted as
an improving step and defaults to 0.1 square model units.

\subsection{Generator Tab}
\label{settings-generator}

The \emph{Generator} tab contains support-generation parameters. Support spacing
sets the nearest-neighbor distance of the hexagonal contact-point lattice and
defaults to 2.0 model units. Boundary contacts use the same value as their
maximum spacing and search radius.

The Contact tip group defines the truncated cones placed at detected contact
points. Top radius defaults to 0.25, bottom radius to 0.75, and height to 2.0
model units. The Generator tab's circumference-points setting controls the
polygonal resolution shared by circular tip and pillar surfaces and defaults
to 12.

The External pillar group controls supports which extend from the build
platform to contact tips. Lattice cell size controls the spatial resolution of
the route search and defaults to 0.5 model units. Minimum support angle is the
smallest permitted angle between a pillar centerline segment and the build
platform and defaults to 30 degrees. The cylindrical platform base defaults to
a height of 0.5 and radius of 2.0 model units. The tapered pillar wire defaults
to a bottom radius of 0.75 and a top radius of 0.5 model units. Model isolation
adds clearance between every external support and the model surface and
defaults to 1.0 model unit.
External pillar bases use the global Z=0 build platform. Models intended for
support generation should therefore be placed in the positive Z half-space;
the \emph{Move to Origin} command performs this placement when appropriate.

The Generator settings control contact placement, tip construction, collision
clearance, route resolution, and surface tessellation. Their use throughout
the support pipeline is described in Chapter~\ref{chap:generate}.

\subsection{Interface Tab}
\label{settings-interface}

The \emph{Interface} tab contains the \emph{Cross-section display distance},
which controls the gap between the selected-model bounding box and the
dedicated cross-section display plane.

Settings persist between runs in \texttt{clip-slicer/settings.ini} beneath the
operating system's user configuration directory.

\section{Mouse Actions}
\label{mouse-actions}

\begin{table}[htbp]
    \centering
    \small
    \begin{tabularx}{\textwidth}{@{}>{\bfseries}p{0.27\textwidth}YY@{}}
        \toprule
        Gesture & Camera target & With Shift: selected-model target \\
        \midrule
        Left drag
            & Orbit the camera around the models.
            & Rotate all selected models around their common geometric center. \\
        Middle drag
            & Translate the camera parallel to the screen.
            & Translate selected models parallel to the screen. \\
        Right drag
            & Move the camera along the axis normal to the screen.
            & Translate selected models along the axis normal to the screen. \\
        Mouse wheel
            & Zoom the camera view in or out.
            & Scale selected models around their common geometric center. \\
        Ctrl+Left drag
            & Duplicate mouse-wheel zoom for a mouse without a wheel.
            & Duplicate mouse-wheel scaling for selected models. \\
        Ctrl+Right drag
            & Duplicate middle-button translation.
            & Duplicate middle-button translation for selected models. \\
        Alt+Mouse wheel
            & \multicolumn{2}{Y@{}}{Move the interactive section plane. This
              action is available while Section mode is enabled.} \\
        Alt+Shift+Left drag
            & \multicolumn{2}{Y@{}}{Duplicate interactive section-plane
              movement for a mouse without a wheel.} \\
        \bottomrule
    \end{tabularx}
    \caption{Mouse controls in the model viewport.}
\end{table}

\section{Keyboard Shortcuts}
\label{keyboard-shortcuts}

\begin{table}[htbp]
    \centering
    \begin{tabularx}{\textwidth}{@{}>{\ttfamily}p{0.24\textwidth}Y@{}}
        \toprule
        \normalfont\bfseries Shortcut & \bfseries Action \\
        \midrule
        Ctrl+O & Open a model in a new document. \\
        Alt+F  & Open the File menu. \\
        Alt+M  & Open the Models menu for visibility and transformation commands. \\
        Alt+S  & Open the Slice menu. \\
        \bottomrule
    \end{tabularx}
    \caption{Application keyboard shortcuts and menu access keys.}
\end{table}

Menu access keys can vary with desktop environment settings. Modifier names in
this manual use the conventional Ctrl, Shift, and Alt labels; the operating
system may display localized equivalents.
